In BGV, BFV and CKKS the ciphertext modulus $Q$ is typically large, much larger than a machine word or register (64 bits on modern computers). Numbers that do not fit inside a register must be handled with multi-precision arithmetic, which is known to be slow due to its serial nature. % that does not allow parallelization. 
High-performance instantiations remedy this by using a \gls{rns} where large numbers are represented as a set of smaller, independent residues via the \gls{crt} (detailed in \autoref{appendix:crt}). \gls{rns} variants of BGV, BFV and CKKS are the standard in practice; OpenFHE, for example, only supports the RNS versions of these schemes~\cite{cryptoeprint:2022/915}.

\begin{definition}[RNS Basis]
    Let $w$ denote the bit-width of a standard machine word. A positive integer $q$ is termed \emph{small} if it fits entirely within this word, such that $0 < q < 2^w$. A set of $k$ moduli $\mathcal{B}_Q = \{q_0, q_1, \dots, q_{k-1}\}$ is considered an \emph{RNS basis} for $Q = \prod_{i=0}^{k-1} q_i$ if its $k$ elements are all small and pairwise coprime ($\gcd(q_i, q_j) = 1$ for all $i \neq j$).
\end{definition}

\begin{definition}[RNS Polynomial]
    \label{def:rns-poly}
    A polynomial $a\in R_Q$ with coefficients $c_j$ is an \emph{RNS polynomial} if it is decomposed into $k$ polynomial \emph{limbs} with coefficients in the basis $\mathcal{B}_Q$ via \gls{crt}, satisfying the explicit ring isomorphism:
    \begin{equation}
        \label{eq:rns-poly}
        R_{Q} \cong R_{q_0} \times R_{q_1} \times \dots \times R_{q_{k-1}}
    \end{equation}

    The limbs are also colloquially termed \emph{towers} or \emph{residues}. We use the notation $[a]_{\mathcal{B}_Q}$ to denote an RNS polynomial $a\in R_Q$ represented entirely by its limbs.
\end{definition}

\newcommand{\firstForm}{\gls{rns} form}
\begin{definition}[First RNS form]
    We define the \emph{\firstForm} of $a$ as the $k \times N$ matrix \eqref{form:crt}, where the $i$-th row contains the coefficients $[c_j]_{q_i}$ of the $i$-th limb.

    \begin{equation}
        \label{form:crt}
        \begin{pmatrix}
            a \bmod q_0 \\
            a \bmod q_1 \\
            \vdots \\
            a \bmod q_{k-1}
        \end{pmatrix}
        =
        \begin{pmatrix}
            c_0 \bmod{q_0} & c_1 \bmod{q_0} & \dots & c_{N-1}\bmod{q_0} \\
            c_0 \bmod{q_1} & c_1 \bmod{q_1} & \dots & c_{N-1}\bmod{q_1} \\
            \vdots & \vdots & \ddots & \vdots \\
            c_0 \bmod{q_{k-1}} & c_1 \bmod{q_{k-1}} & \dots & c_{N-1}\bmod{q_{k-1}}
        \end{pmatrix}
    \end{equation}
\end{definition}

The RNS moduli are specifically chosen to be strictly smaller than $2^w$, so that every coefficient in \eqref{form:crt} is small and fits within a single machine word. Furthermore, the coprimality of the basis ensures this representation is an isomorphism \eqref{eq:rns-poly} allowing reconstruction back to $R_Q$ via the \gls{crt}. Consequently, addition and subtraction of two \gls{rns} polynomials in the same basis is functionally equivalent to the element-wise addition and subtraction of their \firstForm{s} modulo $q_i$. No carries propagate between the towers, allowing all $k$ residues to be processed in parallel utilizing fast, native hardware instructions.

\subsection{Double CRT Polynomial}
While the \firstForm{} perfectly isolates polynomial addition into word-sized registers, standard polynomial multiplication in $R_Q$ operating directly on coefficients scales poorly at $\mathcal{O}(N^2)$. To achieve practical performance, high-speed implementations map these polynomials into an evaluation domain using the \gls{ntt} described in more detail in \autoref{appendix:ntt}.

\begin{definition}[NTT-Friendly RNS Basis]
    \label{def:ntt-friendly}
    Let $N$ denote the dimension of the polynomial ring $R_Q$, where $N$ is a power of two: $N=2^n$. A small modulus $q_i$ is \emph{NTT-friendly} if $q_i \equiv 1 \pmod{2N}$, guaranteeing the existence of a primitive $2N$-th root of unity $\omega_i \in \mathbb{Z}_{q_i}$. An RNS basis $\mathcal{B}_Q$ is \gls{ntt}-friendly if it consists entirely of \gls{ntt}-friendly moduli.
\end{definition}

\newcommand{\secondForm}{evaluation form}
\begin{definition}[Second RNS Form (Evaluation Form)]
    \label{def:ntt-form}
    Let $a \in R_Q$ be an RNS polynomial over an NTT-friendly basis $\mathcal{B}_Q$. The \gls{ntt} is an invertible mapping that transitions a polynomial from its coefficient representation into a point-value representation (see \autoref{appendix:ntt}).
    
    The \emph{\secondForm} (or \emph{NTT form}) of $a$, denoted $\text{NTT}(a)$, is the $k \times N$ matrix obtained by independently applying the $N$-point NTT to each of its $k$ limbs:
    
    \begin{equation}
        \label{form:ntt}
        \Call{NTT}{a}
        \equiv
        \begin{pmatrix}
            \text{NTT}_{q_0}(a \bmod{q_0}) \\
            \text{NTT}_{q_1}(a \bmod{q_1}) \\
            \vdots \\
            \text{NTT}_{q_{k-1}}(a \bmod{q_{k-1}})
        \end{pmatrix}
    \end{equation}
\end{definition}

The \secondForm reduces polynomial multiplication in $R_Q$ to the element-wise product of their respective matrices. The \secondForm{} also inherits the element-wise addition and subtraction from the \firstForm, meaning an entire sequence of arithmetic operations can be executed entirely in parallel using native hardware instructions, without ever leaving the evaluation domain.

%%
%% CLOSING REMARKS 
Although the evaluation form enables fast, parallel polynomial addition and multiplication, certain arithmetic operations fundamentally cannot be performed in the evaluation domain. The transition between representations necessitates a Forward or Inverse NTT, which is a computationally expensive $\mathcal{O}(N \log N)$ procedure. Consequently, a primary objective in the implementation of \gls{rns}-variants of homomorphic encryption schemes is the strategic minimization of transitions between the \gls{rns} and evaluation forms.
